\documentclass[12pt, a4paper]{ctexart}
\usepackage{fontspec}
\usepackage{xcolor}
\usepackage{hyperref}
\usepackage{geometry}
\usepackage{enumitem}
\usepackage{parskip}
\usepackage{titlesec}
\usepackage{tcolorbox}
\tcbuselibrary{breakable}
\usepackage{setspace}
\usepackage{listings}
\usepackage{amssymb}

\geometry{left=2.5cm, right=2.5cm, top=2.5cm, bottom=2.5cm}

\setmainfont{Times New Roman}
\setCJKmainfont{SimSun}

\hypersetup{
    colorlinks=true,
    linkcolor=blue!70!black,
    urlcolor=cyan,
    pdftitle={面向对象下-逐题精讲解义},
    pdfauthor={专升本-fifine老师}
}

\definecolor{myblue}{RGB}{30,100,200}
\definecolor{lightblue}{RGB}{230,240,255}
\definecolor{darkblue}{RGB}{0,50,120}

\newtcolorbox{bluebox}[1][]{
    breakable,
    colback=lightblue,
    colframe=myblue,
    coltitle=darkblue,
    fonttitle=\bfseries,
    title={#1},
    boxrule=0.8pt,
    arc=4pt,
    left=6pt,
    right=6pt,
    top=6pt,
    bottom=6pt,
    before skip=8pt,
    after skip=8pt
}

\usepackage{etoolbox}
\BeforeBeginEnvironment{bluebox}{\par\vfill}%
\AfterEndEnvironment{bluebox}{\par\vfill}%

\titleformat{\section}
  {\normalfont\Large\bfseries\color{myblue}}{\thesection}{1em}{}
\titlespacing*{\section}{0pt}{3.5ex plus 1ex minus .2ex}{1.5ex plus .2ex}

\title{\textcolor{myblue}{\textbf{面向对象（下） - 逐题精讲解义}}}
\author{专升本 fifine老师}
\date{2026年03月02日}

\lstset{
    language=Java,
    basicstyle=\ttfamily\small,
    keywordstyle=\color{blue},
    commentstyle=\color{green!60!black},
    stringstyle=\color{orange},
    numbers=left,
    numberstyle=\tiny\color{gray},
    frame=single,
    breaklines=true,
    columns=fullflexible,
    showstringspaces=false
}

\newcommand{\code}[1]{\texttt{#1}}

\begin{document}

\maketitle
\thispagestyle{empty}
\newpage

\section{面向对象（下）}

\begin{enumerate}[label=\textbf{\arabic*.}]

	\item \textbf{在Java中，以下哪个关键字用于实现继承？}
	      \begin{enumerate}[label=\Alph*.]
		      \item extends
		      \item implements
		      \item interface
		      \item abstract
	      \end{enumerate}

	      \begin{bluebox}{答案与深度解析}
		      \textbf{答案：A. extends}

		      \textbf{【核心认知框架>}
		      \begin{itemize}[leftmargin=*, nosep, label={}]
			      \item \textbf{题目本质}：考查Java中类继承的实现关键字
			      \item \textbf{关键区分点}：extends用于类继承，implements用于接口实现
		      \end{itemize}

		      \textbf{【选项原理深度拆解】}
		      \begin{itemize}[label={}, nosep]
			      \item \textbf{A. extends — 正确（类继承关键字）}
			      \begin{itemize}[label={}, nosep]
				      \item JLS §8.1：使用extends关键字实现类继承
				      \item 子类自动获得父类的非私有成员
				      \item 单继承：一个类只能继承一个父类
			      \end{itemize}
			      \item \textbf{B. implements — 错误（用于接口实现）}
			      \begin{itemize}[label={}, nosep]
				      \item implements用于实现接口
			      \end{itemize}
			      \item \textbf{C. interface — 错误（用于定义接口）}
			      \begin{itemize}[label={}, nosep]
				      \item interface用于定义接口
			      \end{itemize}
			      \item \textbf{D. abstract — 错误（用于定义抽象类/方法）}
			      \begin{itemize}[label={}, nosep]
				      \item abstract用于定义抽象类或抽象方法
			      \end{itemize}
		      \end{itemize}

		      \textbf{【应背诵知识点>}
		      \begin{itemize}[label={}, nosep]
			      \item 类继承：extends（单继承）
			      \item 接口实现：implements（多实现）
			      \item 接口继承：extends（多继承）
		      \end{itemize}
	      \end{bluebox}

	      \vspace{1em}

	\item \textbf{以下关于抽象类的说法错误的是}
	      \begin{enumerate}[label=\Alph*.]
		      \item 抽象类可以有抽象方法
		      \item 抽象类不能被实例化
		      \item 抽象类的子类必须实现所有抽象方法
		      \item 抽象类中可以有非抽象方法
	      \end{enumerate}

	      \begin{bluebox}{答案与深度解析}
		      \textbf{答案：C. 抽象类的子类必须实现所有抽象方法}

		      \textbf{【核心认知框架>}
		      \begin{itemize}[leftmargin=*, nosep, label={}]
			      \item \textbf{题目本质}：考查抽象类的特性与继承规则
			      \item \textbf{关键区分点}：抽象类子类不一定实现所有抽象方法（如果子类也是抽象类）
		      \end{itemize}

		      \textbf{【选项原理深度拆解】}
		      \begin{itemize}[label={}, nosep]
			      \item \textbf{A. 抽象类可以有抽象方法 — 正确}
			      \begin{itemize}[label={}, nosep]
				      \item 抽象方法是抽象类的核心特征
				      \item 只有抽象类才能有抽象方法
			      \end{itemize}
			      \item \textbf{B. 抽象类不能被实例化 — 正确}
			      \begin{itemize}[label={}, nosep]
				      \item 不能使用new创建实例
				      \item 但可以有构造方法（供子类调用）
			      \end{itemize}
			      \item \textbf{C. 抽象类的子类必须实现所有抽象方法 — 错误（本题答案）}
			      \begin{itemize}[label={}, nosep]
				      \item 普通子类：必须实现所有抽象方法
				      \item 抽象子类：可以不实现（继承）
			      \end{itemize}
			      \item \textbf{D. 抽象类中可以有非抽象方法 — 正确}
			      \begin{itemize}[label={}, nosep]
				      \item 抽象类可以有具体实现的方法
			      \end{itemize}
		      \end{itemize}

		      \textbf{【应背诵知识点>}
		      \begin{itemize}[label={}, nosep]
			      \item 抽象类：有抽象方法的类必须声明为abstract
			      \item 抽象方法：只有声明，没有实现
			      \item 子类继承抽象类：普通类需实现，抽象类可继承
		      \end{itemize}
	      \end{bluebox}

	      \vspace{1em}

	\item \textbf{以下哪个关键字用于定义接口？}
	      \begin{enumerate}[label=\Alph*.]
		      \item interface
		      \item abstract class
		      \item class
		      \item extends
	      \end{enumerate}

	      \begin{bluebox}{答案与深度解析}
		      \textbf{答案：A. interface}

		      \textbf{【核心认知框架>}
		      \begin{itemize}[leftmargin=*, nosep, label={}]
			      \item \textbf{题目本质}：考查接口定义关键字
		      \end{itemize}

		      \textbf{【选项原理深度拆解】}
		      \begin{itemize}[label={}, nosep]
			      \item \textbf{A. interface — 正确（定义接口的关键字）}
			      \begin{itemize}[label={}, nosep]
				      \item 接口使用interface关键字定义
			      \end{itemize}
			      \item \textbf{B. abstract class — 错误（定义抽象类）}
			      \item \textbf{C. class — 错误（定义类）}
			      \item \textbf{D. extends — 错误（用于继承）}
		      \end{itemize}

		      \textbf{【应背诵知识点>}
		      \begin{itemize}[label={}, nosep]
			      \item 接口定义：interface 接口名 \{ \}
			      \item JDK8+：接口可以有default方法和static方法
		      \end{itemize}
	      \end{bluebox}

	      \vspace{1em}

	\item \textbf{以下关于接口的说法错误的是}
	      \begin{enumerate}[label=\Alph*.]
		      \item 接口中的方法默认是public abstract修饰的
		      \item 接口中的变量默认是public static final修饰的
		      \item 一个类可以实现多个接口
		      \item 接口可以继承类
	      \end{enumerate}

	      \begin{bluebox}{答案与深度解析}
		      \textbf{答案：D. 接口可以继承类}

		      \textbf{【核心认知框架>}
		      \begin{itemize}[leftmargin=*, nosep, label={}]
			      \item \textbf{题目本质}：考查接口的特性
			      \item \textbf{关键区分点}：接口只能继承接口
		      \end{itemize}

		      \textbf{【选项原理深度拆解】}
		      \begin{itemize}[label={}, nosep]
			      \item \textbf{A. 接口中的方法默认是public abstract修饰的 — 正确}
			      \begin{itemize}[label={}, nosep]
				      \item 方法默认：public abstract
				      \item JDK8+：default方法和static方法除外
			      \end{itemize}
			      \item \textbf{B. 接口中的变量默认是public static final修饰的 — 正确}
			      \begin{itemize}[label={}, nosep]
				      \item 变量默认：public static final（常量）
			      \end{itemize}
			      \item \textbf{C. 一个类可以实现多个接口 — 正确}
			      \begin{itemize}[label={}, nosep]
				      \item 类可以实现多个接口（多实现）
			      \end{itemize}
			      \item \textbf{D. 接口可以继承类 — 错误（本题答案）}
			      \begin{itemize}[label={}, nosep]
				      \item 接口只能继承接口，不能继承类
				      \item 类只能继承类（单继承）
				      \item 类可以实现接口（多实现）
			      \end{itemize}
		      \end{itemize}

		      \textbf{【应背诵知识点>}
		      \begin{itemize}[label={}, nosep]
			      \item 接口：完全抽象的类型
			      \item 接口成员：抽象方法（默认方法、静态方法）、常量
			      \item 类关系：extends（继承类）、implements（实现接口）
			      \item 接口关系：extends（继承接口，可多继承）
		      \end{itemize}
	      \end{bluebox}

	      \vspace{1em}

	\item \textbf{以下哪个关键字用于修饰内部类，使其可以访问外部类的私有成员？}
	      \begin{enumerate}[label=\Alph*.]
		      \item static
		      \item final
		      \item private
		      \item public
	      \end{enumerate}

	      \begin{bluebox}{答案与深度解析}
		      \textbf{答案：A. static}

		      \textbf{【核心认知框架>}
		      \begin{itemize}[leftmargin=*, nosep, label={}]
			      \item \textbf{题目本质}：考查内部类的访问特性
		      \end{itemize}

		      \textbf{【选项原理深度拆解】}
		      \begin{itemize}[label={}, nosep]
			      \item \textbf{A. static — 正确（静态内部类）}
			      \begin{itemize}[label={}, nosep]
				      \item 静态内部类可以访问外部类的静态成员
				      \item 非静态内部类可以访问外部类的所有成员（包括私有）
			      \end{itemize}
			      \item \textbf{B. final — 错误（无此功能）}
			      \item \textbf{C. private — 错误（修饰符不是这样用）}
			      \item \textbf{D. public — 错误（无此功能）}
		      \end{itemize}

		      \textbf{【应背诵知识点>}
		      \begin{itemize}[label={}, nosep]
			      \item 成员内部类：需要外部类对象创建
			      \item 静态内部类：不需要外部类对象
			      \item 局部内部类：定义在方法内
			      \item 匿名内部类：没有类名
		      \end{itemize}
	      \end{bluebox}

	      \vspace{1em}

	\item \textbf{以下关于匿名内部类的说法错误的是}
	      \begin{enumerate}[label=\Alph*.]
		      \item 没有类名
		      \item 可以继承其他类或实现接口
		      \item 只能使用一次
		      \item 不能定义构造方法
	      \end{enumerate}

	      \begin{bluebox}{答案与深度解析}
		      \textbf{答案：C. 只能使用一次}

		      \textbf{【核心认知框架>}
		      \begin{itemize}[leftmargin=*, nosep, label={}]
			      \item \textbf{题目本质}：考查匿名内部类的特性
		      \end{itemize}

		      \textbf{【选项原理深度拆解】}
		      \begin{itemize}[label={}, nosep]
			      \item \textbf{A. 没有类名 — 正确}
			      \item \textbf{B. 可以继承其他类或实现接口 — 正确}
			      \item \textbf{C. 只能使用一次 — 错误（本题答案）}
			      \begin{itemize}[label={}, nosep]
				      \item 每次new都会创建新的匿名内部类
				      \item 可以多次使用
			      \end{itemize}
			      \item \textbf{D. 不能定义构造方法 — 正确}
			      \begin{itemize}[label={}, nosep]
				      \item 匿名内部类没有显式类名
				      \item 无法定义构造方法
			      \end{itemize}
		      \end{itemize}

		      \textbf{【应背诵知识点>}
		      \begin{itemize}[label={}, nosep]
			      \item 语法：new 接口/类() \{ 实现 \}
			      \item 用途：事件处理、线程创建等
			      \item 限制：无类名、无构造方法、只能继承一个类或实现一个接口
		      \end{itemize}
	      \end{bluebox}

	      \vspace{1em}

	\item \textbf{在Java中，以下哪个方法用于比较两个对象是否相等？}
	      \begin{enumerate}[label=\Alph*.]
		      \item equals()
		      \item compareTo()
		      \item isEqual()
		      \item same()
	      \end{enumerate}

	      \begin{bluebox}{答案与深度解析}
		      \textbf{答案：A. equals()}

		      \textbf{【核心认知框架>}
		      \begin{itemize}[leftmargin=*, nosep, label={}]
			      \item \textbf{题目本质}：考查对象相等性比较方法
		      \end{itemize}

		      \textbf{【选项原理深度拆解】}
		      \begin{itemize}[label={}, nosep]
			      \item \textbf{A. equals() — 正确（对象相等比较方法）}
			      \begin{itemize}[label={}, nosep]
				      \item equals()定义在Object类中
				      \item 默认实现：return this == obj（比较引用）
				      \item String类重写了equals()，比较内容
			      \end{itemize}
			      \item \textbf{B. compareTo() — 错误（用于比较顺序）}
			      \begin{itemize}[label={}, nosep]
				      \item 实现Comparable接口的方法
				      \item 返回负数、零、正数
			      \end{itemize}
			      \item \textbf{C. isEqual() — 错误（无此方法）}
			      \item \textbf{D. same() — 错误（无此方法）}
		      \end{itemize}

		      \textbf{【应背诵知识点>}
		      \begin{itemize}[label={}, nosep]
			      \item ==：比较引用地址（基本类型比较值）
			      \item equals()：比较对象内容（需重写）
			      \item hashCode()：与equals()配套使用
		      \end{itemize}
	      \end{bluebox}

	      \vspace{1em}

	\item \textbf{以下关于Java中对象克隆的说法错误的是}
	      \begin{enumerate}[label=\Alph*.]
		      \item 可以通过实现Cloneable接口来实现对象克隆
		      \item clone()方法是在Object类中定义的
		      \item 克隆分为浅克隆和深克隆
		      \item 默认的clone()方法实现的是深克隆
	      \end{enumerate}

	      \begin{bluebox}{答案与深度解析}
		      \textbf{答案：D. 默认的clone()方法实现的是深克隆}

		      \textbf{【核心认知框架>}
		      \begin{itemize}[leftmargin=*, nosep, label={}]
			      \item \textbf{题目本质}：考查对象克隆机制
		      \end{itemize}

		      \textbf{【选项原理深度拆解】}
		      \begin{itemize}[label={}, nosep]
			      \item \textbf{A. 可以通过实现Cloneable接口来实现对象克隆 — 正确}
			      \begin{itemize}[label={}, nosep]
				      \item Cloneable是标记接口
				      \item 必须实现才能调用clone()
			      \end{itemize}
			      \item \textbf{B. clone()方法是在Object类中定义的 — 正确}
			      \begin{itemize}[label={}, nosep]
				      \item protected方法，需重写为public
			      \end{itemize}
			      \item \textbf{C. 克隆分为浅克隆和深克隆 — 正确}
			      \begin{itemize}[label={}, nosep]
				      \item 浅克隆：只复制对象本身
				      \item 深克隆：递归复制引用对象
			      \end{itemize}
			      \item \textbf{D. 默认的clone()方法实现的是深克隆 — 错误（本题答案）}
			      \begin{itemize}[label={}, nosep]
				      \item Object的默认clone()是浅克隆
				      \item 只复制对象本身，引用类型成员共享同一引用
			      \end{itemize}
		      \end{itemize}

		      \textbf{【应背诵知识点>}
		      \begin{itemize}[label={}, nosep]
			      \item 克隆步骤：实现Cloneable + 重写clone() + 改为public
			      \item 浅克隆：基本类型复制值，引用类型复制引用
			      \item 深克隆：所有引用都递归复制
			      \item 序列化方式也可实现深克隆
		      \end{itemize}
	      \end{bluebox}

\end{enumerate}

\end{document}
